\chapter{Introduction}\label{ch:introduction}
% Background, problem statement, purpose and scope of the project.

% ------------------------------------------------------------
% EXAMPLE of inserting a DrawIO diagram:
% 1. Save/edit the diagram in  assets/drawio/<name>.drawio
% 2. Reference it with \drawiofig - the PDF is generated
%    AUTOMATICALLY by latexmk when you compile
%    (requires draw.io Desktop locally).
%
% Usage:          \drawiofig{file-name}{Caption}{fig:label}
% Optional width: \drawiofig[0.5\linewidth]{...}{...}{...}
% NOTE: all three arguments are required!
% ------------------------------------------------------------
\drawiofig{v-model}{Example: DrawIO diagram inserted automatically (the V-model)}{fig:vmodel-example}

As shown in \autoref{fig:vmodel-example}, a diagram is inserted simply by
referencing the name of the .drawio file.

